% Generated draft: requires mathematical and source review.
\documentclass[11pt]{article}
\usepackage[margin=25mm]{geometry}
\usepackage{fontspec}
\setmainfont{Latin Modern Roman}
\usepackage{amsmath,amssymb,mathtools,unicode-math}
\setmathfont{Latin Modern Math}
\usepackage{enumitem,needspace,xurl,hyperref}
\setlength{\parindent}{0pt}
\setlength{\parskip}{5pt}
\setlength{\emergencystretch}{3em}
\newcommand{\N}{\mathbb{N}}
\newcommand{\Z}{\mathbb{Z}}
\newcommand{\Q}{\mathbb{Q}}
\newcommand{\R}{\mathbb{R}}
\newcommand{\C}{\mathbb{C}}
\newcommand{\F}{\mathbb{F}}
\newcommand{\A}{\mathbb{A}}
\newcommand{\PP}{\mathbb{P}}
\newcommand{\E}{\mathbb{E}}
\newcommand{\eps}{\varepsilon}
\newcommand{\dd}{\,\mathrm{d}}
\begin{document}
\section*{UnsolvedMath ID: 30003072}
\textbf{Source problem code:} catalogue entry
\subsection{Definitions and mathematical statement}
For a monoid \(M\) with operation \(\circ\) and identity \(e\), a subset \(F\subseteq M\) generates \(M\) if every element of \(M\) is a finite product of elements of \(F\), allowing the empty product \(e\). Thus \(M\) is finitely generated if
\[
\exists F\subseteq M\quad |F|<\infty,\qquad
M=\bigl\{f_1\circ\cdots\circ f_k : k\in\Z_{\geq 0},\ f_1,\ldots,f_k\in F\bigr\}.
\]
Here \(\circ\) is neutral notation for the monoid operation, not an additional assertion about the construction.

The supplied question is the following: for every birational sequence \(S\) and every total order \(\prec\) that is proper in the intended PBW/Newton--Okounkov construction, is the associated monoid \(\Gamma(S,\prec)\) finitely generated? Equivalently, does every such pair \((S,\prec)\) admit a finite subset \(F\subseteq\Gamma(S,\prec)\) generating that monoid?

\emph{Unresolved scope.} The supplied record does not define a birational sequence, the condition that an order is proper, or the construction of \(\Gamma\). It also does not specify the ambient algebraic group or Lie algebra, the representation-theoretic hypotheses, or the domains over which these data vary. Consequently, the universal quantifiers above retain the source's intended scope; the record alone does not determine a fully self-contained mathematical formulation.
\subsection{Short English statement}
The question asks whether every monoid obtained from a birational sequence and a proper total order in the intended PBW/Newton--Okounkov construction has a finite generating set. The construction and its admissibility conditions are unspecified in the supplied record and require source clarification.
\subsection{Sources}
\begin{itemize}
\item Local UnsolvedMath registry, supplied record with permanent ID \(30003072\), title \emph{Finite Generation of Birational-Sequence Monoids}, source problem code \texttt{catalogue entry}, in the Oberwolfach Reports Open Problems collection.
\item Supplied primary provenance: \emph{Mini-Workshop: PBW Structures in Representation Theory} (2016), Oberwolfach Reports, pp.~615--651, DOI \texttt{10.4171/owr/2016/13}. \url{https://doi.org/10.4171/owr/2016/13}. This is also the supplied literature reference labelled ``OWR report 13/2016.'' The record gives proposed year 2016 and names no individual proposer.
\item Supplied literature reference labelled ``Fang--Fourier--Littelmann'': \url{https://arxiv.org/abs/1510.02295}.
\item Supplied literature reference labelled ``Feigin survey'': \url{https://doi.org/10.4171/ICM2022/81}.
\item Literature assessment dated exactly \texttt{2026-08-21}: the supplied curation labels the problem open, reports finite generation for important reduced-decomposition and classical PBW cases, and states that the report conjectures finite generation for every birational sequence and that no general proof was found. This is supplied dated curation derived from the OWR material, not an independently verified literature assessment; its time-sensitive status and cited claims require independent verification.
\end{itemize}
\end{document}
